\chapter{Introduction}
\section{Problem Definition}
Banks hold detailed records of every account transaction, but turning those records into a decision about a loan application is a different problem. A bank wants to prioritise manual review of applications that are likely to default, so that reviewers spend their time where it matters. The target is binary: a loan labelled \texttt{0} was completed or is running without payment problems, while a loan labelled \texttt{1} was completed unpaid or is running in debt.

The modelling rule is stricter than the prediction task itself. No transaction, balance, or account event dated on or after the loan date may be used to build a predictive feature. A model that violates this rule would score well in development and fail in practice, because it would be reading information that did not exist at decision time. This leakage constraint shapes the warehouse design, the feature mart, and every evaluation reported in this document.

\section{Project Objectives}
The objective is an end-to-end analytics system over the PKDD'99 Czech banking dataset: 682 loans, 76 of them defaults (11.14\%), and 1,056,320 transactions from 1993 to 1998. Specifically, the project aims to:

\begin{enumerate}
    \item Design a Kimball-style galaxy schema that stores loan and transaction facts against conformed date, branch, customer, account, and loan-type dimensions.
    \item Deploy the warehouse in Snowflake with a RAW staging layer and a MART layer of modelled tables and reporting views.
    \item Build a leakage-safe feature mart in which every behavioural feature uses only transactions strictly before the loan date.
    \item Compare six supervised classifiers under repeated cross-validation and select a review model with an explicit metric rule.
    \item Mine association rules and behavioural segments that can be read alongside the model score.
    \item Deliver the results through a Streamlit dashboard that supports review, not just reporting.
\end{enumerate}

\section{Scope and Data Boundary}
The source is a historical, anonymised dataset. It contains no loan product category and no interest rate, so the warehouse keeps a single explicit ``Unspecified'' loan-type member and a null rate rather than inventing values. Branch labels are derived from Czech district names. The dataset is small by modern credit-risk standards, and this project treats the result as a historical demonstration of a method, not as a deployable credit policy.

\section{Report Organization}
Chapter~2 reviews the dataset and the relevant background. Chapter~3 describes the system architecture, Chapter~4 the Snowflake deployment, and Chapter~5 the warehouse design. Chapter~6 covers the supervised learning dataset and leakage control, Chapter~7 pattern mining, Chapter~8 classification, and Chapter~9 clustering. Chapter~10 presents the dashboard, and Chapter~11 concludes. Appendix~A maps each end-semester jury rubric item to its section and dashboard page.
